\section{Surface de travail}

\begin{figure}[H]
\centering
\includegraphics[width=0.8\textwidth]{images/surface de travail.png}
\caption{surface de travail}
\label{fig:surface}
\end{figure}

Pour calculer l'aire de la zone de travail en forme de lentille, la forumle est la suivante ;
\begin{equation}
   A = R^{2}\left(\theta - \sin\theta\right)
\end{equation}

L'angle de travail, dans notre cas, est égal à 

\begin{equation}
   2 \cdot \arccos\left(\frac{227{,}8}{2 \times 215}\right) = 116 \circ
\end{equation}

Ainsi, la surface de travail effectuable par le robot est de 49'493 mm{2}.

